\subsection{Архитектура миссии: конечный автомат}
\label{sec:missions-architecture}

\begin{infobox}[Три принципа]
\begin{enumerate}
    \item \textbf{События, а не время.} Следующая команда --- по событию о завершении предыдущей. Таймер --- только для пауз, которые нужны по смыслу.
    \item \textbf{Ничего не блокировать.} Ни \verb|sleep|, ни циклов ожидания: \verb|callback| должен возвращаться мгновенно.
    \item \textbf{Failsafe.} Удар (\verb|Ev.SHOCK|) --- выключить моторы; разряд (\verb|Ev.LOW_VOLTAGE2|) --- прекратить маршрут и сесть.
\end{enumerate}
\end{infobox}

Когда у миссии несколько этапов, одной переменной-счётчика мало: одно и то же событие \verb|POINT_REACHED| в состоянии «выхожу на окружность» и в состоянии «возвращаюсь в центр» означает разное. Поэтому программа хранит \textbf{состояние} и проверяет его вместе с событием. В светофоре (глава~\ref{sec:leds}) событие было одно --- тик таймера, --- и правила переходов умещались в таблицу \verb|nextState[state]|. В миссии событий несколько, и переход зависит от пары «состояние + событие»; поэтому правила записываются условиями вида \verb|event == Ev.X and state == "Y"|.

\subsubsection*{Окружность по точкам с индикацией}

\begin{codefile}{lua/06_missions/missions_lesson12_circle_points.lua}
\end{codefile}

\paragraph{Что происходит.} Состояния: \verb|PREPARE| $\to$ \verb|FLIGHT| $\to$ \verb|LANDING|. Цвет линейки показывает этап: красный --- программа загружена, белый --- предстарт, зелёный --- маршрут, синий --- посадка, выключено --- на земле. \verb|POINT_REACHED| обрабатывается только в состоянии \verb|FLIGHT|, поэтому случайное событие на посадке ничего не сломает; так же \verb|TAKEOFF_COMPLETE| принимается только в \verb|PREPARE|, а \verb|COPTER_LANDED| гасит линейку только после штатной посадки. \verb|SHOCK| в любом состоянии --- красный, \verb|ENGINES_DISARM| и переход в \verb|EMERGENCY|: из него маршрут уже не продолжается, и красный цвет остаётся.

\subsubsection*{Непрерывная окружность с курсом}

\begin{codefile}{lua/06_missions/missions_lesson14_circle.lua}
\end{codefile}

\paragraph{Что происходит.} Пять состояний: \verb|PREPARE| $\to$ \verb|TO_CIRCLE| (жёлтый, выход в точку $(r, 0)$) $\to$ \verb|CIRCLE| (фиолетовый, таймер двигает цель и курс 12~с) $\to$ \verb|RETURN| (синий, точка над центром) $\to$ \verb|LANDING|. Цепочка \verb|elseif| в \verb|callback| --- это и есть таблица переходов: каждая строка --- «событие и состояние $\to$ действие». \verb|SHOCK| переводит автомат в \verb|EMERGENCY|; таймер окружности и отложенный конец круга (\verb|Timer.callLater(12, ...)|) первым делом проверяют, что автомат всё ещё в \verb|CIRCLE|, поэтому после аварии они не отправят точку и не перекрасят линейку. Посадка в \verb|RETURN| --- только когда \verb|near(0, 0, height)| подтверждает, что дрон над центром (раздел~\ref{sec:missions-circle}). Схема переходов --- на рисунке~\ref{fig:lua-fsm-circle}.

\begin{tikzfig}
\begin{adjustbox}{max width=\linewidth}\begin{tikzpicture}[fig]
\node[fgrey, text width=30mm] (prep) at (0,0) {\textbf{\texttt{PREPARE}}\\ \footnotesize линейка белая};
\node[forange, text width=30mm] (to) at (6.2,0) {\textbf{\texttt{TO\_CIRCLE}}\\ \footnotesize линейка жёлтая};
\node[forange, text width=34mm] (circ) at (12.4,0) {\textbf{\texttt{CIRCLE}}\\ \footnotesize линейка фиолетовая};
\node[fblue, text width=34mm] (ret) at (12.4,-3.2) {\textbf{\texttt{RETURN}}\\ \footnotesize линейка синяя};
\node[fgreen, text width=30mm] (land) at (6.2,-3.2) {\textbf{\texttt{LANDING}}\\ \footnotesize посадка};
\node[fred, text width=30mm] (em) at (0,-3.2) {\textbf{\texttt{EMERGENCY}}\\ \footnotesize красная, моторы выключены};
\draw[farrow] (prep) -- node[flab, above=1pt, font=\ttfamily\scriptsize] {TAKEOFF\_COMPLETE} (to);
\draw[farrow] (to) -- node[flab, above=1pt, font=\ttfamily\scriptsize] {POINT\_REACHED} (circ);
\draw[farrow] (circ) -- node[flab, right=2pt, align=left] {таймер \\ 12~с} (ret);
\draw[farrow] (ret) -- node[flab, above=1pt, align=center, font=\ttfamily\scriptsize] {POINT\_REACHED \\ и near(0, 0, h)} (land);
\draw[dashed, Danger, line width=0.8pt, ->] (prep) -- node[flab, right=2pt, align=left] {\texttt{SHOCK} из любого \\ состояния} (em);
\node[fnote, anchor=north] at (land.south) {\texttt{COPTER\_LANDED}: линейка гаснет};
\end{tikzpicture}\end{adjustbox}
\captionof{figure}{Автомат миссии «окружность»: пять состояний и события, которые переводят автомат из одного в другое. Удар (\texttt{SHOCK}) переводит в \texttt{EMERGENCY} из любого состояния.}\label{fig:lua-fsm-circle}
\end{tikzfig}


\subsubsection*{Аварийные реакции}

\begin{codefile}{lua/06_missions/missions_failsafe.lua}
\end{codefile}

\paragraph{Что происходит.} Дрон патрулирует между двумя точками бесконечно. Состояния: \verb|TAKEOFF| $\to$ \verb|PATROL| $\to$ \verb|LANDING|, при ударе --- \verb|EMERGENCY|. Одно событие получает разную реакцию в зависимости от состояния:
\begin{itemize}
    \item \verb|SHOCK| --- в любом состоянии сразу выключение моторов: аварию не откладывают.
    \item \verb|LOW_VOLTAGE2| в \verb|PATROL| --- посадка на месте. Во время взлёта автопилот команду посадки не примет, поэтому в \verb|TAKEOFF| программа только запоминает разряд (\verb|lowBattery|) и садится сразу после \verb|TAKEOFF_COMPLETE|, не начиная маршрут. В \verb|LANDING| и \verb|EMERGENCY| повторное событие (напряжение может «прыгать» около порога) ничего не меняет.
    \item \verb|POINT_REACHED| ведёт маршрут только в \verb|PATROL|.
\end{itemize}
В симуляторе аварийные события приходят только «по-настоящему»: \verb|SHOCK| --- при столкновении (например, с деревом или стеной), \verb|LOW_VOLTAGE1| и \verb|LOW_VOLTAGE2| --- при разряде батареи в долгом полёте. Кнопки, которая вызывает их вручную, нет. Чтобы проверить реакцию, вызовите обработчик сами, как в главе~\ref{sec:leds}: \texttt{Timer.callLater(12, function() callback(Ev.LOW\_VOLTAGE2) end)}.

\subsubsection*{Запуск с пульта}

На настоящем дроне миссию удобно запускать тумблером: программа загружена, дрон стоит, пилот отходит и переводит тумблер. \verb|Sensors.rc()| возвращает восемь каналов (значения от $-1$ до $1$), тумблер SwA --- восьмой:

\begin{codefile}{lua/06_missions/missions_rc_initiation.lua}
\end{codefile}

\paragraph{Что происходит.} Таймер раз в секунду читает каналы: \verb|table.pack(Sensors.rc())| собирает восемь значений в таблицу, а \verb|rc_chans[8]| --- восьмой из них. Как только он положителен --- флаг \verb|started|, таймер останавливается (иначе он отдавал бы предстарт каждую секунду), предстарт, через 2~с взлёт. Дальше маршрут ведёт уже знакомый счётчик шагов; \verb|Timer.callLater(1, nextPoint)| --- секундное зависание в каждой точке: следующая команда по-прежнему отдаётся только после события, просто с паузой. Настройка пульта и соответствие тумблеров каналам --- в приложении~\ref{sec:rc-control}. В симуляторе роль пульта играет геймпад, подключённый к компьютеру и настроенный в панели \ui{Пульт}; без него все каналы равны нулю, и программа будет ждать тумблер бесконечно. Нет геймпада --- проверьте остальную программу, временно заменив условие запуска на \verb|true|.

\subsubsection*{Светодиоды для отладки}

В примерах этого раздела с индикацией цвет меняется вместе с состоянием. Это дешёвый и надёжный журнал: на полигоне без терминала по цвету видно, на каком этапе программа зависла. Приучите себя подсвечивать состояния с самого начала.

\begin{taskbox}[Автомат с запасом прочности]
\begin{enumerate}
  \item В \texttt{missions\_failsafe.lua} добавьте вызов \verb|callback(Ev.LOW_VOLTAGE2)| через 12~с после запуска. Затем перенесите его на 1.5~с --- пока дрон ещё взлетает.
  \item Добавьте реакцию на \verb|Ev.LOW_VOLTAGE1|: маршрут продолжается, но линейка становится жёлтой, а в журнал печатается \verb|"низкий заряд, этап: "| и текущее состояние.
  \item Добавьте в патруль счётчик кругов: после трёх полных проходов туда и обратно --- штатная посадка, в журнале \verb|"остановка: маршрут пройден"|. Для аварийной посадки печатайте \verb|"остановка: разряд"|.
\end{enumerate}
\expected{1) при 12~с дрон прекращает патруль и садится на месте, линейка красная; при 1.5~с он сначала завершает взлёт, маршрут не начинает и сразу садится; 2) жёлтая линейка и строка с состоянием \texttt{PATROL}, патруль не прерывается; 3) ровно шесть точек патруля, посадка и строка с причиной остановки --- по журналу видно, штатной она была или аварийной.}
\mission{Lua}{71}
\end{taskbox}
